\section{Actual valuative multiplier-ideal contributions}

\begin{definition}[The actual contribution of a divisorial DVR]\label{def:sc0-dvr-multiplier-contribution}
\lean{EnvelopeProperty.Char0.dvrMultiplierContribution}
\leanok
\uses{def:sc0-original-source-dvr-jacobian-ideal,cor:sc0-original-source-dvr-jacobian-finite-order}
Let $C\subseteq V\subseteq K$ be an actual divisorial DVR over its actual
source function field and $\mathfrak a\subset C$ a nonzero actual ideal.
For a rational exponent $q$, put $i=\operatorname{ord}_V(\mathfrak a V)$ and
$j=\operatorname{ord}_V\operatorname{Jac}_V(C)$. Both are finite integers.
Its actual valuative contribution is the inverse image in $C$ of
\[
 \mathfrak m_V^{\max\{0,\lfloor qi\rfloor-j\}}.
\]
This is an actual ideal, with log discrepancy $A_C(V)=1+j$.
\end{definition}

\begin{lemma}[The contribution is exactly the strict discrepancy condition]\label{lem:sc0-dvr-multiplier-contribution-membership}
\lean{EnvelopeProperty.Char0.mem_dvrMultiplierContribution_iff_strict}
\leanok
\uses{def:sc0-dvr-multiplier-contribution,lem:sc0-dvr-ideal-order-principal-product}
An original source element $f$ belongs to that actual contribution exactly
when $f=0$ or
\[
 q\operatorname{ord}_V(\mathfrak a)
 <\operatorname{ord}_V(f)+1+\operatorname{ord}_V\operatorname{Jac}_V(C).
\]
\end{lemma}
\begin{proof}
In an actual DVR, membership in $\mathfrak m^N$ is equivalent to normalized
order at least $N$, by expressing $\mathfrak m=(\pi)$ and using the actual
uniformizer divisibility criterion. Zero belongs to every power. For nonzero
$f$, all three orders are finite integers. The exponent condition is
$\max\{0,\lfloor qi\rfloor-j\}\le\operatorname{ord}_V(f)$, equivalently
$\lfloor qi\rfloor\le\operatorname{ord}_V(f)+j$. An integer $r$ satisfies
$\lfloor qi\rfloor\le r$ exactly when $qi<r+1$. This is the displayed
strict discrepancy inequality. For negative rational exponents the
contribution is the unit ideal; the stated criterion remains valid.
\end{proof}

\begin{lemma}[The original ideal lies in its exponent-one contribution]\label{lem:sc0-original-ideal-multiplier-contribution}
\lean{EnvelopeProperty.Char0.le_dvrMultiplierContribution_one}
\leanok
\uses{def:sc0-dvr-multiplier-contribution}
Every original element of $\mathfrak a$ belongs to its actual exponent-one
valuative contribution.
\end{lemma}
\begin{proof}
Its actual image lies in $\mathfrak a V$, so its order is at least the
actual ideal order $i$. The cutoff exponent at one is
$\max\{0,i-j\}\le i$. Thus its image lies in that actual DVR ideal power,
which is the asserted inverse-image membership.
\end{proof}
